% ===========================================================================
\section{Conclusion}
\label{sec:conclusion}
% ===========================================================================

Pakistan's road transport---the conditional road residual under the current claimant
assignment---emitted \SI{40.19}{\mega\tonne} \COe{} in FY2023--24 and its mixed off-road
machinery \SI{1.73}{\mega\tonne}, together \SI{9.4}{\percent} of a ten-sector national
inventory in which all ten sectors now report the same twelve months. Both are
anchored on the road customer class of the national sectoral sales table, plus metered
transport CNG, so the sector total agrees with Pakistan's fuel statistics by construction
rather than by adjustment.

Three results generalise beyond this country.

\textbf{Anchoring puts the weakest input where it does least damage.} A fleet model that
distributes a known quantity of fuel cannot move the sector total; one that sizes it can. The
class composition reported here is no better supported than any bottom-up equivalent---worse,
in a country with no identified vehicle-kilometre survey---but it is confined to the output
where its weakness matters least. The arrangement is what the IPCC guidance implies where fuel
statistics are strong and fleet data weak, and it is rarely stated as a design choice.

\textbf{A customer-class table is a claimant ledger waiting to be read.} Pakistan's
sectoral sales statistics partition national fuel exhaustively---road, rail, navigation,
power, industry, government---and reading the partition into a two-tier ledger replaced
most of a boundary argument with observation: bulk claimants excluded by the source's own
structure, one carve-out relayed from the sector that measures it, one partition, an
independent cross-check where two unrelated sources report the same boundary row, and the
genuinely unresolved remainder---the retail channel's end-use split---carried as priced
status rows, because naming a residual's direction and sensitivity is more useful than a
correction nobody can evidence.

\textbf{A benchmark match is not provenance.} The previous edition's control total was
arithmetically inconsistent with its own source note, and a matched-year EDGAR comparison
failed to notice, because the wrong number landed near the right vintage. Correcting the
provenance moved the sector six megatonnes and flipped the comparison's direction; on a
scope inferred from EDGAR's category-resolved source files and a fiscal-year-blended
benchmark on the reference period the estimate now sits at \num{0.85} of the blend---a
scope-approximate diagnostic, with the rc4-era \num{0.70} retained in the comparison
table as a superseded mixed-vintage construction. The
paper retains the superseded control as a named executable scenario---the bridge readers
of the earlier edition are owed---and draws the general moral: external agreement tests
consistency, and only provenance tests provenance. The parallel-table \NtwoO{} trap
(Section~\ref{sec:factors}) makes the same point from the factor side: a
\COtwo{}-dominated validation cannot detect a category running \SI{8}{\percent} low
through the wrong trace-gas table.

Fuel controls constrain the GHG magnitude, while end-use claimant evidence constrains the road/off-road boundary. Fleet technologies and operation remain consequential for pollutant factors even when class shares have little effect on fuel-derived CO$_2$. The annual national core covers July 2015--June 2025; the detailed reference-year class analysis does not constitute observed historical fleet composition.

Archived input tables, source-factor registers, build commands, plotted data and numerical checks accompany the sector estimates. The release distinguishes regeneration from compiled tables from a full primary-source extraction.
